\documentclass{article}
\usepackage{graphicx} % Required for inserting images
\usepackage{float}
\usepackage{wrapfig}
\usepackage{caption}
\usepackage[hidelinks]{hyperref}
\begin{document}
\begin{titlepage}

\begin{center}

% Institution logo
\includegraphics[width=0.22\textwidth]{logo.jpg}

\vspace{1cm}

{\Large \textbf{National Initiative on Undergraduate Science (NIUS) Astronomy}}\\[0.3cm]

{\large Batch 21}\\[1cm]

\hrule
\vspace{0.5cm}

{\Huge \textbf{High energy polarization signatures from AGN Jets}}

\vspace{0.3cm}
\hrule

\vspace{2cm}

{\Large \textbf{Project Report}}

\vspace{1cm}

\begin{flushleft}
\large
\textbf{Submitted by:}\\
Dhanad Jha\\
Kirori Mal College, Delhi University

\vspace{0.4cm}

\textbf{Under the Supervision of:}\\
Dr. Bhargav Vaidya\\
Associate Professor\\
Discipline of Astronomy, Astrophysics and Space Engineering\\
Indian Institute of Technology Indore

\vspace{0.4cm}

\end{flushleft}

\end{center}
\end{titlepage}
\newpage

\tableofcontents
\setcounter{tocdepth}{3}
\newpage
\newpage

\section*{Preface}
\addcontentsline{toc}{section}{Preface}
This work attempts to bridge the gap between introductory synchrotron theory and modern AGN polarization research literature. This report aims to create an accessible pathway in AGN polarization physics for beginners. The report presented is a result of multiple literature reviews and computational reproduction of many results discussed throughout the study. 

This project began as part of the "National Initiative on Undergraduate Science (NIUS)" Astronomy (Batch 21) hosted by the Homi Bhabha Centre for Science Education (HBCSE), a national center of the Tata Institute of Fundamental Research (TIFR), Mumbai. This project is being carried out under the supervision of Dr. Bhargav Vaidya, Associate Professor at the Discipline of Astronomy, Astrophysics, and Space Engineering at the Indian Institute of Technology Indore. A part of the work presented in this report was also presented at the Ashoka Physics Students Symposium 3.0, Ashoka University, where it was awarded the First Prize in the poster presentation category.

This report discusses high-energy polarization signatures from AGN jets. The concepts involved in understanding the subject are often fragmented across multiple texts and sources. This report aims to provide an introduction to a wide range of topics, including AGN jets, synchrotron radiation, relativistic effects, polarization, magnetic field orientations, and their physical effects. Understanding the complex phenomena associated with high energy synchrotron emission requires studying these topics and later combining them to interpret the observed results and develop an understanding of the underlying theories. 

Many significant results from various texts have also been reproduced through numerical codes. The GitHub repository containing these codes and figures can be found in the "Code Availability" section at the end of this report.

Working on this project has presented me with multiple learning opportunities. As my first experience in research, it served as an ideal introduction to understanding how scientific research is conducted and the way formal academic studies are carried out. This project helped me bridge the gap between learning theoretical concepts and applying them to understand real physical phenomena in astrophysics.
\newpage

\section*{Acknowledgement}
\addcontentsline{toc}{section}{Acknowledgement}
I would like to thank the organizing team of the National Initiative on Undergraduate Science (NIUS) Astronomy (Batch 21) hosted by the Homi Bhabha Centre for Science Education (HBCSE), a national center of the Tata Institute of Fundamental Research (TIFR), Mumbai, for providing an opportunity to work on this project. A sincere thanks to Prof. (H) Anwesh Mazumdar, who served as academic coordinator for this project.

A sincere thank you to the project supervisor, respected Dr. Bhargav Vaidya, Associate Professor at the Discipline of Astronomy, Astrophysics, and Space Engineering at the Indian Institute of Technology Indore. Without Dr. Vaidya, this project would not have been possible. From the beginning, Dr. Vaidya has guided me at each phase of the study, suggesting materials for the pre-reading phase of the project, structuring the whole work, and guiding at each step of the project with his valuable suggestions and ideas. It has been an incredible learning opportunity working with Dr. Vaidya for the project by interacting with him, developing a deeper understanding of the subject, and learning from his scientific approach and perspective. I am grateful to have had this opportunity to learn from one of the leading experts in the field.

I would also like to thank Mr. Aditya Sharma, PhD Scholar at the Indian Institute of Technology, Indore. Mr. Sharma has immensely helped throughout the project in terms of subject-related discussions, understanding complex topics, and helping develop a clearer understanding of the subject. He has also been instrumental in non-academic support by helping organize accommodation at the institute and other day-to-day activities. Mr. Sharma's support made the project experience memorable while ensuring a better learning experience.

I would also like to thank my parents, Mrs. Sudhira Jha and Mr. Kedar Nath Jha, who have always encouraged and supported me in working towards my goals.I would also like to thank a dear friend, Abhinav Shrivastav, who always acted as a pillar of support, with whom many project-related discussions took place, many of which contributed to ideas discussed in this report.



\newpage
\begin{abstract}

Polarization studies of Active Galactic Nuclei (AGN) jets provide important insight into the magnetic field structure, relativistic dynamics, and radiation mechanisms operating inside some of the most energetic astrophysical systems in the universe. Since synchrotron radiation emitted by relativistic charged particles is naturally polarized, observations of the polarization fraction and the electric vector position angle (EVPA) serve as powerful probes of jet physics. However, interpreting these observations remains challenging due to the combined effects of relativistic motion, Doppler boosting, magnetic field geometry, and particle transport processes.

This report presents a theoretical study of synchrotron polarization in relativistic AGN jets. The report first discusses the foundations of synchrotron radiation and synchrotron emissivity in randomly oriented magnetic fields. Relativistic synchrotron emission recipes used for calculating surface brightness and Stokes parameters from relativistic magnetohydrodynamic (RMHD) simulations are then introduced. The report further explores helical and toroidal magnetic field models proposed to explain large EVPA rotations, polarization swings, and high polarization fractions observed in AGN jets and relativistic outflows. In addition, polarization signatures from different high-energy emission mechanisms and jet geometries are discussed. Several figures and results from the original research papers were reproduced during the project and analyzed in detail. Overall, this report provides an introductory theoretical framework for understanding polarization phenomena in relativistic astrophysical jets.

\vspace{0.3cm}

\noindent \textbf{Keywords:} AGN jets, synchrotron radiation, polarization, Stokes parameters, relativistic jets, magnetic fields, EVPA rotation, RMHD simulations, helical magnetic fields, toroidal magnetic fields

\end{abstract}
\section{Introduction}

The study of polarization in relativistic astrophysical jets has become one of the most important tools for understanding the physical conditions inside Active Galactic Nuclei (AGN). Observations of AGN jets reveal highly energetic plasma outflows capable of emitting radiation across the entire electromagnetic spectrum. A significant fraction of this radiation is synchrotron emission produced by relativistic charged particles spiraling around magnetic field lines. Since synchrotron radiation is naturally polarized, polarization observations provide direct information about the magnetic field structure, particle acceleration mechanisms, and relativistic motion inside the jet.

However, interpreting polarization observations is highly challenging due to the combined effects of relativistic beaming, Doppler boosting, magnetic field geometry, turbulence, and line-of-sight integration. Observations frequently show complex phenomena such as large Electric Vector Position Angle (EVPA) rotations, sudden polarization flips, and variations in polarization fraction that cannot be fully explained using simple random magnetic field models. Understanding these behaviors requires a combination of synchrotron radiation theory, relativistic transformations, and polarization calculations through Stokes parameters.

This report presents a theoretical study of synchrotron polarization in AGN jets. The report begins with the theoretical foundations of synchrotron radiation and polarization, followed by the calculation of synchrotron emissivity and Stokes parameters. The report then discusses different magnetic field configurations including random, helical, and toroidal magnetic field models proposed to explain observed polarization behavior in AGN jets and relativistic outflows. Several reproduced figures and results from important research papers are also discussed throughout the report.

\subsection{Active Galactic Nuclei Jets}

Active Galactic Nuclei are compact regions located at the centers of galaxies and powered by accretion of matter onto supermassive black holes. In many AGN systems, a fraction of the accreting matter is redirected into narrow relativistic plasma outflows called AGN jets. These jets can extend over distances much larger than the host galaxy itself while carrying enormous amounts of energy.

The exact mechanism responsible for launching AGN jets is still an active area of research, but magnetic fields are believed to play a central role. Rotation of the black hole and accretion disk twists magnetic field lines, accelerating plasma outward along the rotation axis. The jets move relativistically with bulk Lorentz factors

\begin{equation}
\Gamma = \frac{1}{\sqrt{1-\beta^2}}
\end{equation}

where:
\begin{itemize}
    \item \(\beta = v/c\)
    \item \(v\) is the jet velocity
    \item \(c\) is the speed of light
\end{itemize}

Because the jet velocity approaches the speed of light, relativistic effects such as Doppler boosting and relativistic aberration strongly modify the observed radiation properties.

AGN jets emit radiation through several physical processes including synchrotron radiation and inverse Compton scattering. Synchrotron radiation dominates the radio through optical bands and often extends into X-ray energies. Since synchrotron radiation depends strongly on magnetic field geometry, polarization observations provide valuable information about the internal magnetic structure of relativistic jets.

\subsection{Polarization of Electromagnetic Radiation}

Polarization describes the orientation and evolution of the electric field vector of an electromagnetic wave. For an electromagnetic wave propagating in the \(x\)-direction, the electric field may oscillate in different directions within the plane perpendicular to propagation.

A general electromagnetic wave may be written as

\begin{equation}
\mathbf{E}(x,t) =
E_y \cos(kx-\omega t)\hat{y}
+
E_z \cos(kx-\omega t+\delta)\hat{z}
\end{equation}

where:
\begin{itemize}
    \item \(E_y\) and \(E_z\) are the field amplitudes
    \item \(k\) is the wave number
    \item \(\omega\) is the angular frequency
    \item \(\delta\) is the phase difference
\end{itemize}

Depending on the value of \(\delta\), different polarization states are obtained:
\begin{itemize}
    \item Linear polarization
    \item Circular polarization
    \item Elliptical polarization
\end{itemize}

In synchrotron radiation, relativistic electrons spiral around magnetic field lines and emit strongly linearly polarized radiation. The polarization direction is directly related to the projected magnetic field geometry on the plane of the sky. Consequently, polarization measurements allow astronomers to infer magnetic field ordering and orientation inside AGN jets.

The polarization fraction is defined as the ratio of polarized intensity to total intensity:

\begin{equation}
\Pi = \frac{I_{\mathrm{pol}}}{I}
\end{equation}

where:
\begin{itemize}
    \item \(I_{\mathrm{pol}}\) is polarized intensity
    \item \(I\) is total intensity
\end{itemize}

Highly ordered magnetic fields generally produce stronger polarization fractions, while random magnetic fields lead to cancellation effects and lower net polarization.

\subsection{Stokes Parameters}

The polarization state of electromagnetic radiation is commonly described using four Stokes parameters: \(I\), \(Q\), \(U\), and \(V\).

The total intensity is

\begin{equation}
I = \langle E_x^2 \rangle + \langle E_y^2 \rangle
\end{equation}

The linear polarization parameters are

\begin{equation}
Q = \langle E_x^2 \rangle - \langle E_y^2 \rangle
\end{equation}

\begin{equation}
U = 2\langle E_xE_y\cos\delta \rangle
\end{equation}

The circular polarization parameter is

\begin{equation}
V = 2\langle E_xE_y\sin\delta \rangle
\end{equation}

where angular brackets denote time averaging.

The degree of linear polarization is calculated through

\begin{equation}
\Pi =
\frac{\sqrt{Q^2+U^2}}{I}
\end{equation}

The observed Electric Vector Position Angle (EVPA) is obtained from

\begin{equation}
\chi =
\frac{1}{2}
\tan^{-1}\left(\frac{U}{Q}\right)
\end{equation}

The Stokes parameter formalism is extremely important in astrophysical polarization studies because it provides a direct mathematical framework for converting theoretical synchrotron emission models into observable polarization quantities.

\subsection{Synchrotron Radiation}

Synchrotron radiation is emitted when relativistic charged particles accelerate in magnetic fields. As relativistic electrons spiral around magnetic field lines, they undergo centripetal acceleration and continuously emit electromagnetic radiation.

The Lorentz force governing particle motion is

\begin{equation}
\mathbf{F} = q(\mathbf{v}\times\mathbf{B})
\end{equation}

where:
\begin{itemize}
    \item \(q\) is the particle charge
    \item \(\mathbf{v}\) is particle velocity
    \item \(\mathbf{B}\) is the magnetic field
\end{itemize}

The characteristic synchrotron frequency for a relativistic electron is

\begin{equation}
\nu_c =
\frac{3eB\sin\alpha}{4\pi mc}\gamma^2
\end{equation}

where:
\begin{itemize}
    \item \(e\) is the electron charge
    \item \(m\) is the electron mass
    \item \(\gamma\) is the electron Lorentz factor
    \item \(\alpha\) is the pitch angle
\end{itemize}

The radiation emitted by an accelerating charge is described using retarded potentials. Since electromagnetic signals propagate at the speed of light, the observed electromagnetic field depends on the position and motion of the particle at an earlier retarded time rather than its instantaneous position. This introduces important relativistic corrections to the observed radiation pattern.

For relativistic particles with \(\gamma \gg 1\), synchrotron radiation becomes strongly beamed along the particle velocity direction within an angular region approximately given by

\begin{equation}
\theta \sim \frac{1}{\gamma}
\end{equation}

This relativistic beaming causes the observed radiation to become highly anisotropic and significantly enhances the observed intensity through Doppler boosting effects. Because synchrotron radiation depends strongly on magnetic field orientation and relativistic motion, it forms the theoretical foundation for modern polarization studies of AGN jets and relativistic astrophysical outflows.
\section{Theoretical Foundations of Synchrotron Emission}
\subsection{Synchrotron Radiation in Random Magnetic Fields}

The work by Crusius and Schlickeiser (1986) presents a detailed theoretical calculation of synchrotron radiation emitted by relativistic charged particles moving in randomly oriented magnetic fields. The primary motivation of this study is to develop a more realistic formalism for astrophysical synchrotron emission, since magnetic fields in many astrophysical environments, such as active galactic nuclei (AGN), supernova remnants, and relativistic jets, are often turbulent or randomly distributed rather than perfectly uniform.

Classical synchrotron radiation theory generally assumes a uniform magnetic field and a fixed pitch angle between the particle velocity and magnetic field direction. However, realistic astrophysical plasmas contain disordered magnetic fields, requiring an averaging procedure over all possible magnetic field orientations and particle pitch angles. The paper derives compact expressions for synchrotron emissivity under such conditions, making the results highly useful for numerical modeling of non-thermal astrophysical sources.

For the synchrotron power emitted by a relativistic electron spiraling in a magnetic field, the characteristic synchrotron frequency is given by

\[
\nu_c = \frac{3eB\sin\alpha}{4\pi mc}\gamma^2 ,
\]

where:
\begin{itemize}
    \item \(e\) is the electron charge,
    \item \(B\) is the magnetic field strength,
    \item \(\alpha\) is the pitch angle,
    \item \(m\) is the electron mass,
    \item \(c\) is the speed of light,
    \item \(\gamma\) is the Lorentz factor of the electron.
\end{itemize}

This equation defines the characteristic frequency at which most synchrotron radiation is emitted by a relativistic electron.

The synchrotron emissivity depends on the synchrotron kernel function,

\[
F(x)=x\int_x^\infty K_{5/3}(\xi)\, d\xi ,
\]

where:
\begin{itemize}
    \item \(x=\nu/\nu_c\),
    \item \(K_{5/3}\) is the modified Bessel function of the second kind.
\end{itemize}

The appearance of the Bessel function originates from the relativistic motion of the charged particle in curved trajectories. 

The major contribution of the paper is the averaging of synchrotron emissivity over isotropically distributed magnetic field orientations. Instead of calculating emission for a single pitch angle, it is integrated over all possible directions of the magnetic field. This produces an angle-averaged emissivity suitable for turbulent astrophysical plasmas.

The particle population is assumed to follow a power-law energy distribution,

\[
N(E)\propto E^{-p},
\]

where:
\begin{itemize}
    \item \(N(E)\) is the number density of particles,
    \item \(E\) is the particle energy,
    \item \(p\) is the spectral index.
\end{itemize}

Using this distribution, the total synchrotron emissivity is obtained by integrating the single-particle emissivity over all electron energies and magnetic field orientations.

This procedure significantly simplifies synchrotron calculations in realistic astrophysical systems. Instead of explicitly tracking complicated magnetic field geometries, the formalism provides compact averaged expressions that can be directly implemented in numerical radiation models.

The formalism can be widely used in studies of AGN jets, blazars, pulsar wind nebulae, and synchrotron self-Compton models. In particular, it provides the theoretical basis for calculating synchrotron spectra in systems where magnetic turbulence plays an essential role.

\subsection{Computing Synchrotron Surface Brightness and Stokes Parameters in AGN Jets}

Synthetic synchrotron emission recipes provide a practical framework for converting relativistic magnetohydrodynamic (RMHD) simulation outputs into observable quantities such as surface brightness, polarization maps, and Stokes parameters. The procedure described by Del Zanna et al. uses physical quantities generated from simulations, including pressure, velocity, density and magnetic field, to calculate the radiation emitted by relativistic charged particles moving in magnetic fields.

The main purpose of this method is to reproduce observational properties of astrophysical jets, especially Active Galactic Nuclei (AGN) jets, by incorporating synchrotron emission processes together with relativistic corrections such as Doppler boosting and polarization angle swing. 

\subsubsection{Calculation of Synchrotron Emission Coefficient}

The local synchrotron emission coefficient $j'_\nu$ in the fluid comoving frame is obtained by integrating the synchrotron power emitted by individual particles over the particle energy distribution:

\begin{equation}
j'_{\nu}(\nu',n')=\int P(\nu',\gamma)f(\gamma)d\gamma
\end{equation}

where:

\begin{itemize}
\item $j'_{\nu}$ = synchrotron emissivity
\item $\nu'$ = observing frequency in the fluid frame
\item $n'$ = observer direction in the fluid frame
\item $P(\nu',\gamma)$ = synchrotron power emitted by one particle
\item $f(\gamma)$ = particle energy distribution function
\item $\gamma$ = particle Lorentz factor
\end{itemize}

The synchrotron power of a single relativistic particle is

\begin{equation}
P(\nu',\gamma)
=
\frac{2e^4}{3m^2c^3}
B_{\perp}^{\prime 2}
\gamma^2
S(\nu',\nu_c')
\end{equation}

where:

\begin{itemize}
\item $e$ = electron charge
\item $m$ = electron mass
\item $c$ = speed of light
\item $B'_{\perp}$ = magnetic field component perpendicular to particle motion
\item $S(\nu',\nu'_c)$ = synchrotron spectral density function
\end{itemize}

The critical synchrotron frequency is given by

\begin{equation}
\nu'_c
=
\frac{3e}{4\pi mc}
B'_{\perp}\gamma^2
\end{equation}

The emitting particles are assumed to follow a power-law distribution:

\begin{equation}
f(\gamma)
=
\frac{K_p}{4\pi}
\frac{p}{mc^2}
\gamma^{-(2\alpha+1)}
\end{equation}

where:

\begin{itemize}
\item $K_p$ = proportionality constant
\item $p$ = thermal pressure
\item $\alpha$ = spectral index
\end{itemize}

Substituting the above expressions into the emissivity integral gives:

\begin{equation}
j'_{\nu}
=
Cp|B'\times n'|^{\alpha+1}
\nu'^{-\alpha}
\end{equation}

where $C$ combines all constants independent of spatial variables.

\subsubsection{Relativistic Corrections}

Since the observer frame differs from the fluid frame, relativistic effects must be included through the Doppler factor:

\begin{equation}
D
=
\frac{1}
{\gamma(1-\beta\cdot n)}
\end{equation}

where:

\begin{itemize}
\item $\beta=v/c$
\item $n$ = observer direction
\item $\gamma$ = bulk Lorentz factor of fluid
\end{itemize}

The observed emissivity becomes

\begin{equation}
j_\nu
=
Cp
|B'\times n'|^{\alpha+1}
D^{\alpha+2}
\nu^{-\alpha}
\end{equation}

The synchrotron cutoff frequency is

\begin{equation}
\nu_\infty
=
D
\frac{3e}{4\pi mc}
|B'\times n'|
\gamma_\infty^2
\end{equation}

where $\gamma_\infty$ represents the maximum particle energy.

The transformed observer direction becomes

\begin{equation}
n'
=
D
\left[
n+
\left(
\frac{\gamma^2}{\gamma+1}
\beta\cdot n
-\gamma
\right)\beta
\right]
\end{equation}

The magnetic field in the comoving frame is

\begin{equation}
B'
=
\frac{1}{\gamma}
\left[
B+
\frac{\gamma^2}{\gamma+1}
(\beta\cdot B)\beta
\right]
\end{equation}

Useful simplified expressions are:

\begin{equation}
B'
=
\frac{1}{\gamma}
\sqrt{
B^2+
\gamma^2
(\beta\cdot B)^2
}
\end{equation}

and

\begin{equation}
|B'\times n'|
=
\frac{1}{\gamma}
\sqrt{
B^2
-
D^2(B\cdot n)^2
+
2\gamma D
(B\cdot n)
(\beta\cdot B)
}
\end{equation}

\subsubsection{Surface Brightness and Polarization Maps}

To obtain observable quantities, the emissivity is integrated along the line of sight.

The surface brightness is

\begin{equation}
I_\nu
=
\int_{-\infty}^{\infty}
j_\nu dX
\end{equation}

where:

\begin{itemize}
\item $X$ represents distance along the line of sight
\end{itemize}

The Stokes parameter $Q$ is

\begin{equation}
Q_\nu
=
\frac{\alpha+1}
{\alpha+\frac{5}{3}}
\int
j_\nu
\cos2\chi
dX
\end{equation}

The Stokes parameter $U$ is

\begin{equation}
U_\nu
=
\frac{\alpha+1}
{\alpha+\frac{5}{3}}
\int
j_\nu
\sin2\chi
dX
\end{equation}

where:

\begin{itemize}
\item $\chi$ = polarization angle
\end{itemize}

The polarization fraction is then calculated as

\begin{equation}
\Pi_\nu
=
\frac{
\sqrt{
Q_\nu^2+U_\nu^2
}
}
{I_\nu}
\end{equation}

To determine $\chi$, the auxiliary quantity $q$ is introduced:

\begin{equation}
q
=
B+n\times(\beta\times B)
\end{equation}

Its components in the sky plane are:

\begin{equation}
q_Y
=
(1-\beta_X)B_Y
+
\beta_YB_X
\end{equation}

\begin{equation}
q_Z
=
(1-\beta_X)B_Z
+
\beta_ZB_X
\end{equation}

The polarization angle terms become

\begin{equation}
\cos2\chi
=
\frac{
q_Y^2-q_Z^2
}
{
q_Y^2+q_Z^2
}
\end{equation}

\begin{equation}
\sin2\chi
=
-
\frac{
2q_Yq_Z
}
{
q_Y^2+q_Z^2
}
\end{equation}

This sequence of equations forms a complete recipe for calculating synchrotron brightness and polarization properties from AGN jet simulations. Starting from particle distributions and emissivity calculations, and ending with Stokes parameters and polarization fraction, this framework provides a practical handbook for converting RMHD simulations into directly observable quantities.
\section{The Helical Magnetic Field Models}

The previous sections introduced the theoretical basis of synchrotron radiation and the calculation of Stokes parameters for relativistic jets. However, observations of many blazars and AGN jets reveal polarization behaviour that cannot be fully explained by random magnetic field configurations alone. In particular, several sources exhibit smooth and large rotations of the Electric Vector Position Angle (EVPA), sometimes extending over more than $180^\circ$ or even several complete rotations. These events are frequently accompanied by changes in polarization fraction and flux. The work of Lyutikov proposes that such behaviour can naturally arise from synchrotron emission within a large-scale helical magnetic field structure inside relativistic jets.

\subsection{Introduction}

The helical magnetic field model assumes that the jet magnetic field contains both longitudinal and toroidal components. Rather than existing as a completely random field, the magnetic field wraps around the jet axis while also extending along the direction of propagation. Such a geometry naturally develops in magnetically launched jets where rotation of the central black hole and accretion disk twists the magnetic field lines.

The total magnetic field may therefore be written as

\begin{equation}
B^2 = B_z^2 + B_{\phi}^2
\end{equation}

where:

\begin{itemize}
\item $B_z$ = magnetic field component along the jet axis
\item $B_\phi$ = toroidal magnetic field component
\end{itemize}

The relative strength of these components determines the pitch angle of the helical structure:

\begin{equation}
\tan \psi
=
\frac{B_\phi}{B_z}
\end{equation}

where:

\begin{itemize}
\item $\psi$ = magnetic pitch angle
\end{itemize}

For small values of $\psi$, the field becomes nearly longitudinal, whereas for large values the toroidal component dominates.

The observed synchrotron emission from this magnetic structure is then calculated using the emissivity and Stokes parameter formalism discussed in the previous section. The local emissivity remains dependent on the perpendicular component of the magnetic field,

\begin{equation}
j_\nu
\propto
|B'\times n'|^{\alpha+1}
\end{equation}

while the total observed polarization properties are obtained through line-of-sight integration of the Stokes parameters.

\subsection{The Sudden Polarization Flip Problem}
\begin{wrapfigure}{r}{0.36\textwidth}
    \centering
    \vspace{-10pt}
    \includegraphics[width=0.34\textwidth]{Fig 1 Lyutikov 2017.png}
    \caption{Optical R-band observations of BL Lac showing the temporal evolution of flux density, polarization degree, and EVPA, from Marscher et al. (2008).}
    \label{fig:bl_lac_observations}
    \vspace{-10pt}
\end{wrapfigure}

One of the most interesting observational features discussed by Lyutikov is the sudden polarization flip. Observations of blazars often show abrupt changes in EVPA close to $90^\circ$. Classical synchrotron theory with random magnetic fields has difficulty explaining this behaviour because random fields generally produce fluctuating polarization rather than smooth systematic rotations.

The problem arises because polarization vectors from different parts of the jet can partially cancel each other. As the observed geometry changes due to relativistic motion or changing viewing angle, the dominant contribution to polarization can switch between different magnetic field directions.

This transition causes the observed polarization angle to suddenly rotate by

\begin{equation}
\Delta \chi
\approx
90^\circ
\end{equation}

while simultaneously producing a temporary reduction in polarization fraction.

Physically, this behaviour occurs because two competing polarization directions become nearly equal in magnitude. Small changes in geometry then make one contribution dominate over the other.

\subsubsection{Assumptions and Physical Description of the Model}

The model uses several important assumptions:

\begin{enumerate}
\item The jet moves relativistically with bulk Lorentz factor $\Gamma$.
\item The magnetic field possesses an ordered helical structure.
\item The electron population follows a power-law energy distribution.
\item Synchrotron emission dominates the observed radiation.
\item Relativistic Doppler boosting modifies observed emission.
\end{enumerate}
\begin{figure}[H]
    \centering
    \includegraphics[width=0.45\textwidth]{Lyutikov 2017.png}
    \caption{A schematic relativistic jet model with a helical magnetic field producing polarization variations. Figure 2 from \textit{Polarization Swings in Blazars} by Maxim Lyutikov et al. (2017).}
    \label{fig:blazar_swing_model}
\end{figure}
The Doppler factor remains

\begin{equation}
D
=
\frac{1}
{\Gamma(1-\beta\cos\theta)}
\end{equation}

where:

\begin{itemize}
\item $\theta$ = observer viewing angle
\item $\beta=v/c$
\end{itemize}

Because observed intensity depends strongly on the Doppler factor,

\begin{equation}
I_\nu
\propto
D^{3+\alpha}
\end{equation}

small changes in viewing geometry can substantially modify the observed brightness and polarization properties.

The polarization fraction is still obtained from the Stokes parameters introduced in the previous section:

\begin{equation}
\Pi
=
\frac{\sqrt{Q^2+U^2}}
{I}
\end{equation}

The observed EVPA is calculated from

\begin{equation}
\chi
=
\frac{1}{2}
\tan^{-1}
\left(
\frac{U}{Q}
\right)
\end{equation}

As the helical field rotates relative to the observer, the values of $Q$ and $U$ evolve continuously. Since the polarization angle depends on their ratio, a sign change in either quantity can produce large rotations in EVPA.

Thus the model explains polarization swings as a purely geometrical effect without requiring rapid physical changes in the jet itself.

%\begin{figure}[h]
%\centering
%\includegraphics[width=0.75\textwidth]{figure3.png}
%\caption{Reproduction of Figure 3 from Lyutikov et al. showing polarization evolution associated with helical magnetic field geometry. This figure was reproduced during the project.}
%\end{figure}

%Figure 3 illustrates how changes in geometry affect the observed polarization vectors and produce large EVPA rotations.
\begin{figure}[!ht]
\centering

\includegraphics[width=0.72\textwidth]{Lytikov Fig 5.png}
\caption{Polarization fraction and EVPA for a jet undergoing planar motion with bulk Lorentz factor $\gamma=10$, observed at minimum viewing angles $\theta_{ob,0}=1/(5\gamma)$, $1/\gamma$, and $2/\gamma$ (left to right). Top panels show polarization as a function of oscillation angle for different intrinsic pitch angles, while bottom panels show EVPA (solid line) and the normalized EVPA variation rate $\dot{\phi}_{PA}$ (dashed line).}

\vspace{0.4cm}

\includegraphics[width=0.72\textwidth]{Lytikov Fig 6.png}
\caption{Same set-up as Fig.~5. Top panels show polarization as a function of position angle (PA), while bottom panels show intensity as a function of polarization degree $\Pi$.}

\end{figure}
\subsubsection{Model Performance and Limitations}

The helical magnetic field model successfully reproduced several important observational properties of AGN jets, including smooth EVPA rotations, sudden $90^\circ$ polarization flips, and simultaneous variations in polarization fraction and flux. The model demonstrated that these complex behaviors can emerge naturally from relativistic effects and changing viewing geometry without requiring sudden physical changes inside the jet. Figures 3, 5 and 6 reproduced during this project further support the model predictions. However, the model assumes a largely ordered magnetic field and does not fully account for turbulence, shocks, or magnetic reconnection processes that may exist in real jets. Overall, the model provides a physically consistent framework linking synchrotron theory with observed polarization behavior in AGN jets.
\subsection{Polarization as a Probe of High-Energy Emission Mechanisms}

\subsubsection{Distinguishing Leptonic and Hadronic Emission Models}

Although the helical magnetic field model successfully explained large EVPA rotations and polarization swings, observations of BL Lac objects revealed an additional challenge concerning the origin of high-energy X-ray and $\gamma$-ray emission. Multiple observations showed that different theoretical models could reproduce nearly identical spectral energy distributions while proposing very different physical origins for the radiation. Consequently, spectral fitting alone became insufficient for identifying the true emission mechanism.

Two major scenarios were commonly proposed. The first is the leptonic scenario, where relativistic electrons produce synchrotron emission at lower energies and inverse Compton scattering at higher energies. The second is the hadronic scenario, where relativistic protons contribute significantly to high-energy radiation through proton synchrotron processes and secondary particle cascades.

The motivation of the study by Tavecchio et al. was to determine whether polarization observations could distinguish these competing mechanisms. Since polarization depends strongly on magnetic field geometry and emission processes, the authors proposed that X-ray and $\gamma$-ray polarization measurements could act as an additional diagnostic tool beyond conventional spectral studies.


\subsubsection{Axisymmetric Helical Jet Model}

The model assumes a stationary and axisymmetric relativistic jet possessing an ordered magnetic field structure. The emitting region was represented by a cylindrical jet geometry where the jet radius evolves according to

\begin{equation}
R_0 \propto z_0^q
\end{equation}

where:

\begin{itemize}
\item $R_0$ = jet radius
\item $z_0$ = distance along jet axis
\item $q$ = jet expansion parameter
\end{itemize}

The magnetic field was assumed to contain both longitudinal and toroidal components, consistent with the helical magnetic field framework discussed previously.

The particle population follows a power-law energy distribution:

\begin{equation}
N(\gamma)
=
N_0\gamma^{-p}
\end{equation}

where:

\begin{itemize}
\item $N_0$ = normalization constant
\item $\gamma$ = particle Lorentz factor
\item $p$ = particle spectral index
\end{itemize}

The model further assumes:

\begin{enumerate}
\item Relativistic jet motion
\item Synchrotron dominated emission
\item Ordered magnetic field geometry
\item Power-law particle distribution
\item Axisymmetric jet structure
\end{enumerate}

The synchrotron emissivity was calculated using the formalism introduced previously where the emission depends on the magnetic field component perpendicular to the observer direction:

\begin{equation}
j_\nu
\propto
|B'\times n'|^{\alpha+1}
\end{equation}

This formulation allows the model to connect jet geometry directly with observable radiation properties.


\subsubsection{Polarization and Stokes Parameter Calculations}

The polarization calculations were performed through the Stokes parameter formalism introduced previously. The total intensity is written as

\begin{equation}
I=
\int j_\nu dV
\end{equation}

The linear polarization components become

\begin{equation}
Q=
\int \Pi_0j_\nu \cos(2\chi)dV
\end{equation}

\begin{equation}
U=
\int \Pi_0j_\nu \sin(2\chi)dV
\end{equation}

where:

\begin{itemize}
\item $\Pi_0$ = intrinsic polarization
\item $\chi$ = polarization angle
\item $dV$ = volume element
\end{itemize}

The polarization fraction is then calculated as

\begin{equation}
\Pi
=
\frac{\sqrt{Q^2+U^2}}
{I}
\end{equation}

Using these equations, polarization predictions for different emission scenarios were obtained.



\subsubsection{Predictions and Physical Insights}

The study demonstrated that polarization observations can break degeneracies that remain unresolved through spectral fitting alone. The results showed that leptonic models generally predict weaker high-energy polarization, while hadronic models predict substantially larger polarization fractions.

\begin{figure}[h]
\centering
\includegraphics[width=0.75\textwidth]{Tavecchio Image.jpg}
\caption{Reproduction of Figure 2 from Tavecchio et al. showing X-ray polarization degree ($\Pi_X$, top panels) and chromaticity ($\Pi_O/\Pi_X$, bottom panels) plotted as functions of the proton emission region width $\Delta z$. The left panels assume constant proper number density, while the right panels assume constant magnetization, for a viewing angle $\theta_{\mathrm{obs}} = 0.1$ rad.}
\end{figure}

The reproduced results suggest that polarization measurements can become a powerful tool for identifying the physical origin of high-energy emission in AGN jets. The study further extends the helical magnetic field framework by connecting magnetic field structure and synchrotron radiation directly with competing emission mechanisms.

\subsection{Jet Shape Effects and Non-Thermal Electron Population Models}


Previous studies mainly focused on understanding polarization through magnetic field configurations such as random fields and helical magnetic structures. However, observations of AGN jets frequently show polarization properties that change significantly across different observing frequencies. These variations are often difficult to explain solely through magnetic field geometry. This raised an additional question: how do the physical shape of the jet and the spatial distribution of emitting particles affect observed polarization properties?

The study by Bolis et al. aimed to investigate how jet geometry and particle injection mechanisms influence synchrotron polarization. The authors proposed a model in which the structure of the jet itself can significantly modify polarization signatures. Two simple scenarios for populating the jet with non-thermal electrons were considered. The first assumes particle acceleration throughout the entire jet volume, while the second assumes localized injection regions associated with shocks or magnetic reconnection sites.


\subsubsection{Non-Thermal Relativistic Jet Structure and Particle Population Assumptions}

The jet was modeled as an expanding relativistic outflow where the shape of the jet determines the evolution of both particle density and magnetic field properties. Unlike simplified cylindrical models, the geometry here directly affects the physical conditions along the jet.

The particle population follows a power-law energy distribution:

\begin{equation}
N(\gamma)
=
K\gamma^{-p}
\end{equation}

where:

\begin{itemize}
\item $K$ = normalization constant
\item $\gamma$ = particle Lorentz factor
\item $p$ = spectral index
\end{itemize}

The model assumes:

\begin{enumerate}
\item Relativistic jet motion
\item Synchrotron dominated emission
\item Non-thermal particle populations
\item Ordered magnetic field evolution
\item Frequency-dependent radiation processes
\end{enumerate}

The synchrotron emissivity follows the standard form introduced previously:

\begin{equation}
j_\nu
\propto
|B' \times n'|^{\alpha+1}
\end{equation}

where the observed radiation depends on the magnetic field component perpendicular to the observer direction.

The two electron population scenarios considered are:

\begin{enumerate}
\item Uniform particle injection throughout the jet volume
\item Localized particle injection near acceleration regions
\end{enumerate}

These assumptions allow the model to study how particle distribution and jet expansion together influence the observed polarization.


\subsubsection{Polarization and Stokes Parameter Calculations}

Polarization properties were calculated using the same Stokes parameter formalism introduced previously for synchrotron radiation calculations. The total observed intensity is obtained through line-of-sight integration:

\begin{equation}
I_\nu
=
\int j_\nu dX
\end{equation}

The linear polarization components are written as

\begin{equation}
Q_\nu
=
\int j_\nu \cos(2\chi)dX
\end{equation}

\begin{equation}
U_\nu
=
\int j_\nu \sin(2\chi)dX
\end{equation}

where:

\begin{itemize}
\item $\chi$ = polarization angle
\item $X$ = line-of-sight coordinate
\end{itemize}

The polarization fraction is then determined through

\begin{equation}
\Pi
=
\frac{\sqrt{Q_\nu^2+U_\nu^2}}
{I_\nu}
\end{equation}

Using these quantities, the model predicts how polarization changes with observing frequency and jet structure. Since the radiation is integrated over the entire emitting region, contributions from different parts of the jet may partially cancel each other, producing variations in net polarization.


\subsubsection{Jet Structure Effects and Observational Signatures}

The results showed that the physical shape of the jet has a strong influence on observed polarization properties. Small variations in jet expansion produced significant changes in the contribution from different emitting regions and therefore altered the net polarization.

The study further showed that localized electron injection regions generally produce stronger polarization signatures because the radiation originates from smaller regions with relatively ordered field structures. In contrast, uniform particle distributions tend to produce greater cancellation effects and therefore lower polarization fractions.

The work demonstrated that polarization behaviour depends not only on magnetic field geometry but also on particle injection mechanisms and jet dynamics. The results suggest that realistic AGN jet models must include both magnetic structure and particle transport effects to accurately reproduce multifrequency polarization observations.

\section{Electromagnetically Dominated Toroidal Magnetic Field Model}

\subsection{Introduction and Motivation}

An important development in the study of relativistic polarization phenomena was introduced by Lyutikov, Pariev, and Blandford in their work on polarization from relativistically expanding magnetized outflows. The model was originally proposed to explain the extremely large linear polarization detected in prompt gamma-ray burst (GRB) emission by the RHESSI satellite. Conventional hydrodynamic fireball models struggled to explain such high polarization because magnetic fields generated locally at shocks possess only small coherence lengths and therefore produce strong cancellation effects in the observed polarization.

The authors proposed that the outflow is instead electromagnetically dominated, meaning that large-scale ordered magnetic fields are present from the beginning of the explosion and transport most of the energy outward in the form of Poynting flux. In this framework, the magnetic field retains coherence over angular scales larger than the visible relativistic patch \(R/\Gamma\), naturally producing large observable polarization fractions.

The model assumes that the relativistic ejecta form a quasi-spherical expanding shell carrying a predominantly toroidal magnetic field. During expansion, most of the magnetic energy becomes concentrated inside a thin emitting shell surrounding the outflow. The shell expands relativistically with bulk Lorentz factor \(\Gamma\), while synchrotron radiation is emitted by relativistic electrons accelerated within the shell.

The toroidal magnetic field geometry is represented as

\begin{equation}
\mathbf{B}' = B_{\phi}\{-\sin\phi,\cos\phi,0\}
\end{equation}

where:
\begin{itemize}
    \item \(B_{\phi}\) is the toroidal magnetic field strength
    \item \(\phi\) is the azimuthal angle around the jet axis
\end{itemize}

The observer is located at a viewing angle \(\theta_{\mathrm{ob}}\) relative to the symmetry axis of the shell. Because of relativistic beaming, the observer only receives radiation from a narrow visible region of angular size approximately \(1/\Gamma\). When the viewing angle satisfies \(\theta_{\mathrm{ob}} \gtrsim 1/\Gamma\), the visible region contains an approximately uniform magnetic field orientation, producing large net polarization.

The model further assumes:
\begin{enumerate}
    \item Relativistic quasi-spherical shell expansion
    \item Dominantly toroidal magnetic field geometry
    \item Synchrotron dominated radiation
    \item Isotropic power-law electron distribution
    \item Optically thin emission region
    \item Negligible circular polarization
\end{enumerate}

The relativistic electrons follow a power-law energy distribution:

\begin{equation}
dn = K_e \epsilon^{-p} d\epsilon \, dV \, d\Omega_p
\end{equation}

where:
\begin{itemize}
    \item \(K_e\) is the normalization constant
    \item \(\epsilon\) is the electron energy
    \item \(p\) is the particle spectral index
    \item \(dV\) is the volume element
    \item \(d\Omega_p\) is the momentum-space solid angle
\end{itemize}

The synchrotron emissivity depends on the magnetic field component perpendicular to the observer direction in the comoving frame:

\begin{equation}
j_{\nu} \propto \left| \mathbf{B}' \times \mathbf{n}' \right|^{\alpha + 1}
\end{equation}

where:
\begin{itemize}
    \item \(\mathbf{n}'\) is the photon propagation direction in the comoving frame
    \item \(\alpha\) is the synchrotron spectral index
\end{itemize}

This formulation directly connects the relativistic magnetic field geometry with observable polarization properties.

\subsection{Stokes Parameter Calculation for a Quasi-Spherical Thin Emitting Shell}

The polarization calculations were performed using a relativistically expanding quasi-spherical thin shell model. The shell possesses cylindrical symmetry and carries a globally ordered toroidal magnetic field.

The shell expands radially with velocity

\begin{equation}
\mathbf{v} = \beta \{\sin\theta \cos\phi,\sin\theta \sin\phi,\cos\theta\}
\end{equation}

where:
\begin{itemize}
    \item \(\beta = v/c\)
    \item \(\theta\) and \(\phi\) are spherical polar coordinates
\end{itemize}

The observer direction is represented by

\begin{equation}
\mathbf{n} = \{\sin\theta_{\mathrm{ob}},0,\cos\theta_{\mathrm{ob}}\}
\end{equation}

The Doppler boosting factor becomes

\begin{equation}
D = \frac{1}{\Gamma(1-\mathbf{n}\cdot\mathbf{v})}
\end{equation}

where:
\begin{itemize}
    \item \(\Gamma\) is the bulk Lorentz factor
    \item \(D\) determines relativistic boosting of the emitted radiation
\end{itemize}

Because synchrotron emissivity depends on the local magnetic field orientation in the comoving frame, Lorentz transformations must be applied carefully. The photon direction in the shell frame is

\begin{equation}
\mathbf{n}' =
\frac{
\mathbf{n}
+\Gamma \mathbf{v}
\left(
\frac{\Gamma}{\Gamma+1}\mathbf{n}\cdot\mathbf{v}-1
\right)
}{
\Gamma(1-\mathbf{n}\cdot\mathbf{v})
}
\end{equation}

The angle between the photon direction and magnetic field in the comoving frame is

\begin{equation}
\cos\chi' = \hat{\mathbf{B}}' \cdot \mathbf{n}'
\end{equation}

leading to

\begin{equation}
\sin^2\chi' =
1 -
\frac{
\sin^2\phi \sin^2\theta_{\mathrm{ob}}
}{
\Gamma^2(1-\beta l)^2
}
\end{equation}

where

\begin{equation}
l =
\cos\theta \cos\theta_{\mathrm{ob}}
+
\sin\theta \sin\theta_{\mathrm{ob}} \cos\phi
\end{equation}

The observed polarization cannot be calculated simply by summing intensities because relativistic motion rotates the polarization vector in the observer frame. This relativistic swing of polarization introduces partial depolarization even when the magnetic field remains perfectly ordered.

The pulse-integrated Stokes parameters are therefore written as

\begin{equation}
I =
\int
D^{2+\frac{p-1}{2}}
B'^{\frac{p+1}{2}}
|\sin\chi'|^{\frac{p+1}{2}}
\, d\Omega
\end{equation}

\begin{equation}
Q =
\int
D^{2+\frac{p-1}{2}}
B'^{\frac{p+1}{2}}
|\sin\chi'|^{\frac{p+1}{2}}
\cos(2\tilde{\chi})
\, d\Omega
\end{equation}

\begin{equation}
U =
\int
D^{2+\frac{p-1}{2}}
B'^{\frac{p+1}{2}}
|\sin\chi'|^{\frac{p+1}{2}}
\sin(2\tilde{\chi})
\, d\Omega
\end{equation}

where:
\begin{itemize}
    \item \(\tilde{\chi}\) is the observed polarization position angle
    \item \(d\Omega = \sin\theta d\theta d\phi\)
    \item \(I,Q,U\) are the observed Stokes parameters
\end{itemize}

Due to cylindrical symmetry of the shell,

\begin{equation}
U = 0
\end{equation}

and the polarization fraction becomes

\begin{equation}
\Pi = \frac{|Q|}{I}
\end{equation}

The theoretical maximum synchrotron polarization for a perfectly ordered field is

\begin{equation}
\Pi_{\max} =
\frac{p+1}{p+\frac{7}{3}}
\end{equation}

However, relativistic aberration and integration over different emitting regions reduce the observed polarization below this ideal value.

The model further assumes:
\begin{enumerate}
    \item Instantaneous thin-shell emission
    \item Uniform emissivity across the shell
    \item Optically thin synchrotron radiation
    \item Negligible Faraday rotation
    \item Negligible random magnetic field component
    \item No circular polarization (\(V=0\))
\end{enumerate}

Figure~\ref{fig:lyutikov_fig2} shows the reproduced polarization emissivity map from the toroidal shell model. The figure illustrates how relativistic beaming confines the observed emission into a narrow visible patch while simultaneously producing relativistic rotation of polarization vectors.

\begin{figure}[!htbp]
    \centering
    \includegraphics[width=0.38\textwidth]{Figure 2 Lyutikov 2003.png}
    \caption{Reproduction of Figure 2 from Lyutikov et al. showing polarization emissivity maps for different $\Gamma = 1$. The arrows represent local polarization directions while the intensity distribution illustrates relativistic Doppler boosting effects.}
    \label{fig:lyutikov_fig2}
\end{figure}

\vspace{0.2cm}

Figure~\ref{fig:lyutikov_fig3} presents the predicted dependence of polarization fraction on observer viewing angle for different Lorentz factors.

\begin{figure}[!htbp]
    \centering
    \includegraphics[width=0.65\textwidth]{Figure 3 Lyutikov 2003.png}
    \caption{Reproduction of Figure 3 from Lyutikov et al. showing polarization fraction as a function of observer viewing angle for different Lorentz factors and particle spectral indices.}
    \label{fig:lyutikov_fig3}
\end{figure}
\subsection{Discussion and Physical Interpretation}

The toroidal magnetic field model successfully demonstrated that highly polarized synchrotron emission can naturally arise from electromagnetically dominated relativistic outflows. One of the most important conclusions of the work is that large observed polarization fractions require magnetic field coherence scales larger than the visible relativistic patch \(R/\Gamma\). Such coherence cannot easily be maintained in purely hydrodynamic shock-generated magnetic field models.

The reproduced results show that polarization strongly depends on observer viewing geometry. When the observer lies close to the symmetry axis (\(\theta_{\mathrm{ob}} < 1/\Gamma\)), the visible emitting region contains magnetic field vectors with many different orientations. Their contributions partially cancel, producing relatively weak net polarization. As the observer moves outside the relativistic beaming cone, the visible patch becomes more uniform in magnetic field orientation, causing the polarization fraction to increase substantially.

The model predicts polarization fractions approaching approximately \(50\%-60\%\) for particle spectral indices near \(p=3\). Although synchrotron radiation in a perfectly uniform stationary magnetic field can theoretically produce even larger polarization fractions, relativistic aberration and differential rotation of polarization vectors reduce the observed value. This relativistic depolarization effect represents one of the major physical insights of the study.

The reproduced polarization maps further illustrate the importance of relativistic Doppler boosting. For large Lorentz factors, the observable radiation originates only from a narrow region with angular size approximately \(1/\Gamma\). Consequently, even though the shell itself is quasi-spherical, the observer effectively detects radiation from a small localized emitting patch.

Another important prediction of the model is that the observed electric vector position angle remains approximately fixed during the burst because it is determined primarily by the projection of the symmetry axis onto the observer plane. This prediction differs significantly from many hydrodynamic jet models that often predict rapid EVPA variations associated with jet-break evolution.

The model also established an important connection between polarization and causal structure in relativistic outflows. Magnetically dominated plasmas possess relativistic Alfv\'en velocities capable of re-establishing causal contact across the visible shell region. In contrast, purely hydrodynamic relativistic outflows remain causally disconnected over scales larger than \(1/\Gamma\). This causal reconnection allows the magnetic field to preserve large-scale coherence during expansion.

Despite its success, the model contains several idealizations. The calculations neglect turbulent magnetic fields, magnetic reconnection, and internal plasma propagation effects such as Faraday depolarization. Real relativistic jets are expected to contain both ordered and disordered magnetic field components, which would reduce the net observable polarization. The model also assumes a thin-shell approximation and time-integrated emission rather than performing fully time-dependent radiative transfer calculations.

Nevertheless, the toroidal magnetic field framework provides a physically consistent explanation for large polarization fractions observed in relativistic astrophysical outflows. The work further established polarization as a powerful diagnostic tool for distinguishing between hydrodynamic and electromagnetically dominated jet models. The reproduced figures and calculations performed during this project strongly support the conclusion that large-scale ordered magnetic fields play a major role in relativistic high-energy astrophysical jets.

\section{Conclusion and Discussion}

This report presented a theoretical study of synchrotron polarization in relativistic Active Galactic Nuclei (AGN) jets and explored how magnetic field geometry, relativistic motion, and particle dynamics influence observable polarization properties. Beginning with the basic theory of synchrotron radiation and Stokes parameters, the report developed the mathematical framework required to connect relativistic magnetized plasma models with observable quantities such as surface brightness, polarization fraction, and Electric Vector Position Angle (EVPA).

The study first introduced synchrotron emissivity in random magnetic fields and the relativistic synchrotron emission formalism used in relativistic magnetohydrodynamic (RMHD) simulations. The calculation of Stokes parameters through line-of-sight integration was discussed in detail since these quantities form the foundation of modern polarization studies. The report then explored several important magnetic field models proposed to explain observational polarization behavior in AGN jets.

The helical magnetic field model demonstrated that large EVPA rotations and sudden polarization flips can naturally emerge from relativistic viewing geometry and ordered magnetic structures without requiring rapid intrinsic changes inside the jet. The work further showed how relativistic Doppler boosting and geometric effects strongly modify observed polarization signatures.

The report also examined polarization as a diagnostic tool for distinguishing between leptonic and hadronic high-energy emission models. The reproduced results from Tavecchio et al. showed that hadronic emission scenarios generally predict larger X-ray and \(\gamma\)-ray polarization fractions due to stronger magnetic field ordering in proton synchrotron emission. These results demonstrate that polarization observations can help break degeneracies that remain unresolved through spectral fitting alone.

The influence of jet geometry and non-thermal particle injection mechanisms was further investigated using the models proposed by Bolis et al. The results showed that polarization depends not only on magnetic field structure but also on jet expansion, particle transport, and localized acceleration regions. In particular, localized particle injection regions were found to produce stronger polarization signatures because they reduce cancellation effects from multiple emitting regions.

Finally, the electromagnetically dominated toroidal magnetic field model developed by Lyutikov, Pariev, and Blandford was studied in detail. The reproduced calculations demonstrated that highly polarized synchrotron emission can naturally arise from relativistically expanding magnetized outflows possessing large-scale ordered toroidal magnetic fields. The work highlighted the importance of relativistic aberration, Doppler boosting, and causal magnetic field coherence in producing large observable polarization fractions.

Several figures and tables from the original research papers were successfully reproduced during this project. These reproductions provided practical understanding of the numerical implementation of polarization calculations and strengthened the connection between theoretical models and observational predictions.

Overall, this report demonstrates that polarization serves as one of the most powerful probes of magnetic field structure and relativistic plasma physics in AGN jets. Although many simplified assumptions were used in the discussed models, including ordered magnetic fields and idealized jet geometries, the results collectively show that relativistic polarization studies provide critical insight into the physical processes governing high-energy astrophysical jets. Future studies involving turbulence, magnetic reconnection, time-dependent radiative transfer, and fully three-dimensional RMHD simulations will further improve our understanding of polarization phenomena in relativistic astrophysical outflows.

\section{Code Availability}

The numerical codes, reproduced figures, simulation outputs, and supporting computational material developed during this project are publicly available on GitHub at:

\begin{center}
\href{https://github.com/Dhanad1905/High-Energy-Polarization-Signatures-From-AGN-Jets#high-energy-polarization-signatures-from-agn-jets}{High-Energy-Polarization-Signatures-From-AGN-Jets GitHub Repository}
\end{center}

The repository contains Python implementations of synchrotron emissivity calculations, Stokes parameter computations, polarization mapping routines, and helical magnetic field models discussed throughout this report. The repository also includes reproduced figures, Jupyter notebooks, and documentation related to the numerical methods used in this study.

\newpage

\begin{thebibliography}{99}
\addcontentsline{toc}{section}{References}
\bibitem{AandA2023}
``Polarization Signatures in Relativistic AGN Jets,''
\textit{Astronomy and Astrophysics},
2023.
doi: \texttt{10.1051/0004-6361/202348050}

\bibitem{AandA2024}
``High-Energy Polarization Signatures in Relativistic Jets,''
\textit{Astronomy and Astrophysics},
2024.
doi: \texttt{10.1051/0004-6361/202450387}

\bibitem{Bols2017}
Bols, E., et al.,
``Polarization Signatures from Structured Relativistic Jets,''
\textit{Monthly Notices of the Royal Astronomical Society},
vol. 468, no. 1, pp. 1239--1252, 2017.
doi: \texttt{10.1093/mnras/stx359}

\bibitem{Crusius1986}
Crusius, A., \& Schlickeiser, R.,
``Synchrotron Radiation in Random Magnetic Fields,''
\textit{Astronomy and Astrophysics},
vol. 164, pp. L16--L18, 1986.

\bibitem{DelZanna2006}
Del Zanna, L., Volpi, D., Amato, E., \& Bucciantini, N.,
``Simulated Synchrotron Emission from Pulsar Wind Nebulae,''
\textit{Astronomy and Astrophysics},
vol. 453, no. 2, pp. 621--633, 2006.
doi: \texttt{10.1051/0004-6361:20064858}

\bibitem{Ghatak}
Ghatak, A.,
\textit{Optics},
6th ed.,
McGraw Hill Education, 2017.

\bibitem{Longair}
Longair, M. S.,
\textit{High Energy Astrophysics},
3rd ed.,
Cambridge University Press, 2011.

\bibitem{Lyutikov2003}
Lyutikov, M., Pariev, V. I., \& Blandford, R. D.,
``Polarization of Prompt Gamma-Ray Burst Emission,''
\textit{The Astrophysical Journal},
submitted to arXiv, 2003.
doi: \texttt{10.48550/arXiv.astro-ph/0305410}

\bibitem{PhysicsReports2016}
``Synchrotron Radiation and Polarization in Relativistic Jets,''
\textit{Results in Physics},
vol. 6, pp. 1175--1185, 2016.
doi: \texttt{10.1016/j.rinp.2016.10.021}

\bibitem{PhysRevD2024}
``Relativistic Polarization and Magnetic Field Geometry in AGN Jets,''
\textit{Physical Review D},
vol. 110, 123032, 2024.
doi: \texttt{10.1103/PhysRevD.110.123032}

\bibitem{Rybicki}
Rybicki, G. B., \& Lightman, A. P.,
\textit{Radiative Processes in Astrophysics},
Wiley-VCH, 1979.

\bibitem{Tavecchio2018}
Tavecchio, F., et al.,
``Probing the Emission Mechanism of Blazars through X-ray and Gamma-ray Polarization,''
arXiv preprint, 2018.
doi: \texttt{10.48550/arXiv.1808.08960}

\bibitem{Arxiv2025}
``Recent Developments in Relativistic Jet Polarization Models,''
arXiv preprint, 2025.
doi: \texttt{10.48550/arXiv.2505.19784}

\bibitem{marscher2008}
Marscher, A. P., Jorstad, S. G., D'Arcangelo, F., Smith, P. S., Williams, G. G., Larionov, V. M., Oh, H., Olmstead, A. R., Aller, M. F., Aller, H. D., et al. (2008). ``The inner jet of an active galactic nucleus as revealed by a radio-to-$\gamma$-ray outburst.'' \textit{Nature}, \textbf{452}, 966--969.

\end{thebibliography}
\end{document}